\begin{abstract}
Video pretraining gives an embodied agent a model of how the visual world
evolves. Embodied interaction asks for something stronger, because an agent has
to know how its own actions participate in that evolution. We formulate this as
one modelling problem and call it Joint Action-World Modeling. A single
continuous flow-matching process generates the future observation and the future
action chunk together, the two streams carry independently drawn noise levels,
and they exchange information through one shared attention space at every layer
of the backbone. The pair of noise levels spans a plane, and the training
distribution covers that plane, so one trained model yields three inference
paths: a joint path that produces actions and futures together, an
action-clamped ray that turns the model into an action-conditioned world model,
and a world-clamped ray that turns it into an inverse dynamics model. Those
paths are also the measurement instruments of the paper, because they read out
how much action information the world representation carries and how strongly
the generated future follows an intervention on the action. The formulation
extends one step further to action consequence, defined as the displacement of
world keypoints that an action actually produces in the image, which supplies
grounding on human video that carries no action label and gives a probe of
representation content that does not depend on the action space of any single
embodiment. Coupling becomes a single experimental axis with three degrees of
freedom, namely cross-stream visibility, gradient flow, and coupling depth, and
that axis separates joint modelling from an auxiliary action head under matched
parameters, data, and optimisation. The paper delivers the formulation, the
implementation, the pretraining corpus of robot and egocentric human footage,
and a pre-registered measurement programme with declared floors, criteria, and
failure conditions. Reference values reproduced from a public world-action
implementation fix the performance regime the mainline reproduction has to
reach, and every value that the programme still has to measure is marked as a
target in both the text and the source.
\end{abstract}
